\chapter{Symmetry Counts}

\begin{center}\itshape
Symmetry, as wide or as narrow as you may define its meaning, is one idea by which man through the ages has tried to
comprehend and create order, beauty, and perfection.\\
\upshape --- Hermann Weyl, \emph{Symmetry} (1952)
\end{center}

In 1875 Arthur Cayley, a lawyer turned professor, tried to count the molecules of chemistry. An alkane $\mathrm C_n\mathrm H_{2n+2}$
is a tree of $n$ carbon atoms, each of degree at most four, and two drawings are the same molecule when one can be moved onto
the other. Cayley's numbers for small $n$ matched the chemists' isomers, but his method was cumbersome and some of his larger counts
later proved wrong; the trouble was the symmetry. A molecule with a symmetry is drawn in fewer different ways than
one without, and a count that ignores this overcounts the lopsided ones and undercounts the symmetric ones.

Sixty years later George Pólya, in Zürich, settled the matter with a single idea. Do not count the objects; count, for each
symmetry, the drawings it leaves unchanged, and average. His 1937 paper turned chemistry into algebra and became a classic
of combinatorics. The averaging lemma at its heart is usually called Burnside's lemma, because William
Burnside put it in his 1897 textbook. But Frobenius had published it in 1887, and Cauchy had used it in 1845. When Peter
Neumann pointed this out in 1979, in a paper called ``A lemma that is not Burnside's'', the name was already too famous to
change.

This chapter is about that averaging idea and what grows out of it. Counting orbits is computing an expectation: the mean
number of fixed points of a random symmetry. From there it is a short walk to random permutations, whose cycles behave like
independent Poisson variables, and to characters, which count solutions of equations in groups.

\begin{think}
Shuffle a deck of $n$ cards perfectly. On average, how many cards stay in their place? Now replace the shuffle by a uniformly
random rotation of a cube, and the cards by its $6$ faces, or its $8$ vertices, or its $12$ edges. Can you guess the average
number of fixed faces, vertices and edges before computing anything? What do the three situations have in common?
\end{think}

\section{Counting up to symmetry}

When a finite group $G$ acts on a finite set $X$, Burnside's lemma (Tome~2, Chapter~3) says that the number of orbits is
$\frac1{|G|}\sum_g|X^g|$, the average number of fixed points. The problems below apply it to dice, colourings, words and
graphs, and in one case to a whole infinite family of actions at once.

\begin{exercise}[$\star$]\label{pb:dice}
A die is a cube whose faces carry the numbers $1$ to $6$, each once. Two dice are the same if a rotation carries one onto the
other.
\begin{enumerate}[label=(\alph*)]
\item How many different dice are there?
\item On a standard die, opposite faces add up to $7$. How many different standard dice are there?
\item Show that no die is the same as its mirror image, so that the dice of (a) come in mirror pairs.
\end{enumerate}
\end{exercise}

\begin{exercise}[$\star\star$]\label{pb:patterns}
Let $G$ act on a finite set $X$, and let $K$ be a set of $k$ colours. Two colourings $f:X\to K$ are \emph{equivalent} if
$f'=s\circ f\circ g^{-1}$ for some $g\in G$ and some permutation $s$ of $K$: the symmetry moves the points and the colours may
be renamed.
\begin{enumerate}[label=(\alph*)]
\item Show that for $g\in G$ and $s\in\mathfrak S(K)$, the number of colourings $f$ with $s\circ f\circ g^{-1}=f$ is
\[
\prod_{\text{cycles of }g}\ \Bigl(\sum_{j\mid\ell}j\,c_j(s)\Bigr),
\]
where the product is over the cycles of $g$ on $X$, $\ell$ is the length of the cycle, and $c_j(s)$ is the number of $j$-cycles
of $s$.
\item Count the necklaces of $n$ beads in two colours, up to rotation and up to exchanging the two colours. Check that there
are $8$ for $n=6$.
\item In how many essentially different ways can one paint the faces of a cube with at most three colours, if the colours
themselves do not matter (only which faces share a colour)?
\end{enumerate}
\end{exercise}

\begin{exercise}[$\star\star$]\label{pb:necklacering}
For an integer $a\ge1$, let $M(n,a)$ be the number of primitive necklaces of length $n$ over an alphabet of $a$ letters:
necklaces whose $n$ rotations are pairwise distinct.
\begin{enumerate}[label=(\alph*)]
\item Show that $\sum_{d\mid n}d\,M(d,a)=a^n$. Deduce that $M(n,a)=\frac1n\sum_{d\mid n}\mu(d)a^{n/d}$, and that
$\sum_{d\mid n}\mu(d)a^{n/d}\equiv0\pmod n$ for \emph{every} integer $a$, negative ones included.
\item Prove the \emph{cyclotomic identity}, in $\Z[[x]]$:
\[
\frac1{1-ax}=\prod_{n\ge1}\Bigl(\frac1{1-x^n}\Bigr)^{M(n,a)}.
\]
\item (Metropolis and Rota.) Show that for all $a,b\ge1$,
\[
M(n,ab)=\sum_{\operatorname{lcm}(i,j)=n}\gcd(i,j)\,M(i,a)\,M(j,b),
\]
and explain the identity with necklaces over the alphabet of pairs of letters.
\end{enumerate}
\end{exercise}

\begin{exercise}[$\star\star\star$]\label{pb:unlabelled}
Let $g_n$ be the number of graphs on $n$ vertices up to isomorphism, and $N=\binom n2$.
\begin{enumerate}[label=(\alph*)]
\item Show that $g_n\ge\frac{2^N}{n!}$.
\item Show that $g_n\sim\frac{2^N}{n!}$: almost all graphs have no symmetry.
\item Show more precisely that $\frac{n!\,g_n}{2^N}=1+\frac{n(n-1)}{2^{n-1}}+O\bigl(n^42^{-3n/2}\bigr)$.
\end{enumerate}
\end{exercise}

\begin{exercise}[$\star\star\star$]\label{pb:commuting}
Let $h_k(n)$ be the number of $k$-tuples $(\sigma_1,\dots,\sigma_k)$ of permutations of $\{1,\dots,n\}$ that commute pairwise,
with $h_k(0)=1$. Show that
\[
\sum_{n\ge0}h_2(n)\frac{x^n}{n!}=\prod_{j\ge1}\frac1{1-x^j},\qquad\sum_{n\ge0}h_3(n)\frac{x^n}{n!}=\prod_{j\ge1}\Bigl(\frac1{1-x^j}\Bigr)^{\sigma(j)},
\]
where $\sigma(j)$ is the sum of the divisors of $j$. In particular $h_2(n)=n!\,p(n)$, with $p(n)$ the number of partitions of $n$.
\end{exercise}

\begin{exercise}[$\star\star$]\label{pb:landau}
Let $G$ be a finite group with exactly $k$ conjugacy classes.
\begin{enumerate}[label=(\alph*)]
\item Show that $|G|\le k^{2^{k-1}}$. So, for each $k$, only finitely many groups have $k$ classes.
\item Find all groups with at most $3$ conjugacy classes.
\end{enumerate}
\end{exercise}

\section{What a random symmetry fixes}

Burnside's lemma is a statement about an expectation. Make $g$ uniformly random in $G$ and let $F(g)=|X^g|$. Then
$\mathbb E[F]$ is the number of orbits, and the higher moments count orbits on pairs, triples, and so on. Problem~\ref{pb:fixmoments}
did this for $\mathfrak S_n$; here the group is arbitrary.

\begin{exercise}[$\star\star$]\label{pb:rank}
Let $G$ act transitively on a set $X$ of $n\ge2$ points, let $g$ be uniform in $G$ and $F$ its number of fixed points. The
\emph{rank} $r$ of the action is the number of orbits of $G$ on $X\times X$.
\begin{enumerate}[label=(\alph*)]
\item Show that $\mathbb E[F]=1$ and $\mathbb E[F^2]=r\ge2$.
\item Show that the following are equivalent: $G$ is $2$-transitive (transitive on ordered pairs of distinct points);
$r=2$; $\operatorname{Var}F=1$; the character $\pi-1$ is irreducible, where $\pi(g)=F(g)$ is the permutation character.
\item Find the rank of the dihedral group $D_n$ acting on the $n$ vertices of a regular polygon.
\end{enumerate}
\end{exercise}

\begin{exercise}[$\star\star$]\label{pb:derangements}
Let $G$ act transitively and faithfully on a set $X$ of $n\ge2$ points.
\begin{enumerate}[label=(\alph*)]
\item (Cameron and Cohen.) Show that at least a proportion $\frac1n$ of the elements of $G$ have no fixed point.
\item Show that equality holds if and only if $G$ is \emph{sharply $2$-transitive}: for all pairs $(x,y)$ and $(x',y')$ of
distinct points there is exactly one $g$ with $gx=x'$ and $gy=y'$. Check it for the group of maps $t\mapsto at+b$ of
$\mathbb F_p$, $a\ne0$.
\item (Jordan.) Deduce that a finite group is never the union of the conjugates of a proper subgroup.
\end{enumerate}
\end{exercise}

\begin{exercise}[$\star\star$]\label{pb:glfix}
Let $A$ be a uniformly random matrix in $GL_n(\mathbb F_q)$, $n\ge2$, and $d=\dim\ker(A-I)$.
\begin{enumerate}[label=(\alph*)]
\item Show that $\mathbb E[q^d]=2$ and $\mathbb E[q^{2d}]=q+3$.
\item Deduce that
\[
\frac1q\le P(1\text{ is an eigenvalue of }A)\le\frac1{q-1}.
\]
\end{enumerate}
\end{exercise}

\begin{exercise}[$\star\star\star$]\label{pb:ksubsets}
Let $\pi_k(\sigma)$ be the number of $k$-element subsets of $\{1,\dots,n\}$ fixed by $\sigma\in\mathfrak S_n$, and let
$0\le k,l\le\frac n2$.
\begin{enumerate}[label=(\alph*)]
\item Show that $\langle\pi_k,\pi_l\rangle=\min(k,l)+1$.
\item Deduce that $\chi_0=1$ and $\chi_k=\pi_k-\pi_{k-1}$ ($1\le k\le\frac n2$) are distinct irreducible characters of
$\mathfrak S_n$, of degrees $\binom nk-\binom n{k-1}$.
\item (Livingstone and Wagner.) Let $G\le\mathfrak S_n$. Show that the number of orbits of $G$ on $k$-subsets is a
non-decreasing function of $k$ for $0\le k\le\frac n2$.
\end{enumerate}
\end{exercise}

\section{Random permutations}

A uniformly random permutation is the simplest random symmetry, and its cycles are a world of their own. Tome~3 met them
through Feller's coupling: the number of cycles is a sum of independent Bernoulli variables. Here we look at which cycles
there are.

\begin{exercise}[$\star$]\label{pb:samecycle}
Let $\sigma$ be a uniformly random permutation of $\{1,\dots,n\}$.
\begin{enumerate}[label=(\alph*)]
\item Show that the length of the cycle containing $1$ is uniform on $\{1,\dots,n\}$.
\item Show that $1,2,\dots,k$ lie in the same cycle with probability $\frac1k$, for $1\le k\le n$.
\end{enumerate}
\end{exercise}

\begin{exercise}[$\star$]\label{pb:foata}
Write a permutation $\sigma$ in cycle notation with each cycle starting with its largest element, and the cycles in increasing
order of their first elements; then erase the parentheses. For example $(3\,1)(5\,2\,4)\mapsto3\,1\,5\,2\,4$.
\begin{enumerate}[label=(\alph*)]
\item Show that this is a bijection from $\mathfrak S_n$ to the words that use each of $1,\dots,n$ once, and that the number of
cycles of $\sigma$ is the number of records of the word (letters larger than all the letters before them).
\item Deduce that the number of records of a random word and the number of cycles of a random permutation have the same
law, and that exactly $\stirlingi nk$ words have $k$ records.
\end{enumerate}
\end{exercise}

\begin{exercise}[$\star\star$]\label{pb:roots}
Let $a_n(k)$ be the number of $\sigma\in\mathfrak S_n$ with $\sigma^k=e$, with $a_0(k)=1$.
\begin{enumerate}[label=(\alph*)]
\item Show that $\sum_na_n(k)\frac{x^n}{n!}=\exp\bigl(\sum_{d\mid k}\frac{x^d}d\bigr)$.
\item Let $p$ be prime. Show that $a_n(p)=a_{n-1}(p)+\frac{(n-1)!}{(n-p)!}a_{n-p}(p)$ for $n\ge p$, and deduce that $p$ divides
$a_n(p)$ for every $n\ge p$.
\item Show that $a_n(p)\equiv p\pmod{p(p-1)}$ for $n\ge p$.
\end{enumerate}
\end{exercise}

\begin{exercise}[$\star\star$]\label{pb:oddcycles}
Let $o_n$ and $e_n$ be the numbers of permutations of $\{1,\dots,n\}$ whose cycles all have odd length, respectively all have
even length.
\begin{enumerate}[label=(\alph*)]
\item Show that $o_n=e_n=\bigl(1\cdot3\cdot5\cdots(n-1)\bigr)^2$ for $n$ even, and $o_n=n\,o_{n-1}$ for $n$ odd.
\item Deduce that the probability that a random permutation has odd order is asymptotic to $\sqrt{\frac2{\pi n}}$.
\end{enumerate}
\end{exercise}

\begin{exercise}[$\star\star\star$]\label{pb:squares}
Let $s_n$ be the number of permutations of $\{1,\dots,n\}$ that are squares, $\sigma=\tau^2$ for some $\tau\in\mathfrak S_n$.
\begin{enumerate}[label=(\alph*)]
\item Show that $\sigma$ is a square if and only if, for every even $\ell$, $\sigma$ has an even number of cycles of length $\ell$.
\item Show that
\[
\sum_{n\ge0}s_n\frac{x^n}{n!}=\sqrt{\frac{1+x}{1-x}}\ \prod_{m\ge1}\cosh\frac{x^{2m}}{2m}.
\]
\item Deduce that $s_{2m+1}=(2m+1)s_{2m}$, and that the probability that a random permutation is a square is at least that of
Problem~\ref{pb:oddcycles}(b).
\end{enumerate}
\end{exercise}

\begin{exercise}[$\star\star$]\label{pb:poissoncycles}
Let $C_j$ be the number of $j$-cycles of a uniformly random permutation of $\{1,\dots,n\}$, and $(c)_m=c(c-1)\cdots(c-m+1)$.
\begin{enumerate}[label=(\alph*)]
\item Show that for all $m_1,\dots,m_k\ge0$,
\[
\mathbb E\Bigl[\prod_{j=1}^k(C_j)_{m_j}\Bigr]=\begin{cases}\prod_jj^{-m_j}&\text{if }\sum_jjm_j\le n,\\0&\text{otherwise.}\end{cases}
\]
\item Deduce that for fixed $k$, $(C_1,\dots,C_k)$ converges in law to independent Poisson variables of parameters
$1,\frac12,\dots,\frac1k$.
\item Show that the probability that a random permutation has no cycle shorter than $m$ tends to
$e^{-H_{m-1}}\approx\frac{e^{-\gamma}}m$ as $n\to\infty$, where $H_{m-1}=1+\frac12+\dots+\frac1{m-1}$. Compare with the
density of the integers that have no prime factor smaller than $m$.
\end{enumerate}
\end{exercise}

\begin{figure}[ht]
\centering
\fig{ch04-cycles}
\caption{The laws of the numbers $C_1$, $C_2$, $C_3$ of fixed points, $2$-cycles and $3$-cycles of a random permutation of
$30$ points, computed exactly (bars), against Poisson laws of parameters $1$, $\frac12$, $\frac13$ (black marks). At this size
the difference is invisible: it is smaller than $10^{-8}$ for $C_1$ (Problem~\ref{pb:poissoncycles}).}
\label{fig:cycles}
\end{figure}

\begin{exercise}[Journey]\label{pb:golomb}
\emph{The longest cycle.} Let $L_n$ be the length of the longest cycle of a uniformly random permutation of
$\{1,\dots,n\}$, and $E_1(u)=\int_u^\infty\frac{e^{-v}}v\,dv$.
\begin{enumerate}[label=(\alph*)]
\item Fix $0<x<1$. Draw $N$ with $P(N=n)=(1-x)x^n$ for $n\ge0$, then a uniform permutation of $\{1,\dots,N\}$, and let
$C_j$ be its number of $j$-cycles. Show that the $C_j$, $j\ge1$, are independent Poisson variables of parameters $\frac{x^j}j$.
\item Deduce that
\[
\sum_{n\ge0}(1-x)x^n\,\mathbb E[L_n]=\sum_{m\ge0}\Bigl(1-\exp\Bigl(-\sum_{j>m}\frac{x^j}j\Bigr)\Bigr).
\]
\item Let $x=e^{-1/t}$. Show that the right-hand side is equivalent to $t\int_0^\infty\bigl(1-e^{-E_1(u)}\bigr)du$ as
$t\to\infty$.
\item Deduce that if $\frac{\mathbb E[L_n]}n$ has a limit $\lambda$, then
\[
\lambda=\int_0^\infty\bigl(1-e^{-E_1(u)}\bigr)du=\int_0^\infty e^{-u-E_1(u)}\,du=0.62433\ldots
\]
\end{enumerate}
\end{exercise}

\begin{figure}[ht]
\centering
\fig{ch04-longest}
\caption{For $n=1,\dots,60$, computed exactly: the mean proportion $\mathbb E[L_n]/n$ of the points that lie in the longest
cycle (blue), which tends to the Golomb--Dickman constant $\lambda=0.6243\ldots$ (Problem~\ref{pb:golomb}), and the
probability that the longest cycle has more than $\frac n2$ points (yellow), which tends to $\log2$ (Tome~3, Exercise~2.6).}
\label{fig:longest}
\end{figure}

\begin{exercise}[$\star\star$]\label{pb:conjugate}
Let $\sigma$ and $\tau$ be independent uniform permutations of $\{1,\dots,n\}$, and $P_n$ the probability that they are
conjugate, with $P_0=1$.
\begin{enumerate}[label=(\alph*)]
\item Show that $\sum_{n\ge0}P_n=\prod_{j\ge1}\sum_{m\ge0}\frac1{j^{2m}m!^2}=4.2634\ldots$
\item Show that $\liminf n^2P_n\ge\sum_{n\ge0}P_n$.
\end{enumerate}
\end{exercise}

\section{Characters and factorisations}

Characters turn counting problems in a group into sums over its irreducible representations. The number of solutions of an
equation such as $xyx^{-1}y^{-1}=g$, or $x^2=g$, or $ab=c$ with $a$, $b$ in given classes, is a short formula in the
character table. For the last problem of the chapter there is also a picture.

\begin{exercise}[$\star\star$]\label{pb:commutators}
Let $G$ be a finite group and $N(g)$ the number of pairs $(x,y)\in G^2$ with $xyx^{-1}y^{-1}=g$.
\begin{enumerate}[label=(\alph*)]
\item (Frobenius, 1896.) Show that $N(g)=|G|\sum_\chi\frac{\chi(g)}{\chi(1)}$, the sum over the irreducible characters.
\item Deduce from the character table of $A_5$ (Tome~2, Exercise~11.7) that every element of $A_5$ is a commutator,
and compute $N(g)$ on each class.
\end{enumerate}
\end{exercise}

\begin{exercise}[$\star\star\star$]\label{pb:frobschur}
Let $G$ be finite, $r(g)$ the number of $x\in G$ with $x^2=g$, and for an irreducible character $\chi$ let
$\nu(\chi)=\frac1{|G|}\sum_{x\in G}\chi(x^2)$.
\begin{enumerate}[label=(\alph*)]
\item Show that $\nu(\chi)\in\{-1,0,1\}$, and that $\nu(\chi)\ne0$ if and only if $\chi$ is real-valued.
\item (Frobenius and Schur, 1906.) Show that $r=\sum_\chi\nu(\chi)\chi$.
\item Show that $\nu=-1$ for the irreducible character of degree $2$ of the quaternion group $Q_8$.
\item Show that $\nu(\chi)=1$ when the representation can be realised by real matrices. Every irreducible representation of
$\mathfrak S_n$ can be (Young, 1900). Deduce that the sum of the degrees of the irreducible characters of $\mathfrak S_n$ is
the number of $\sigma$ with $\sigma^2=e$.
\end{enumerate}
\end{exercise}

\begin{exercise}[$\star\star\star$]\label{pb:ncycles}
Let $n\ge2$ and $c=(1\,2\,\dots\,n)$. Let $V\subset\C^n$ be the hyperplane $\{\sum x_i=0\}$, with $\mathfrak S_n$ permuting
the coordinates, and $\chi_r$ the character of its exterior power $\Lambda^rV$, $0\le r\le n-1$. Take as known that the
$\Lambda^rV$ are irreducible and pairwise non-isomorphic (Fulton and Harris, Proposition~3.12).
\begin{enumerate}[label=(\alph*)]
\item Show that $\chi_r(c)=(-1)^r$, and that every other irreducible character vanishes at $c$.
\item Show that the number of ways to write $g=ab$ with $a$ in a class $A$ and $b$ in a class $B$ is
$\frac{|A||B|}{|G|}\sum_\chi\frac{\chi(a)\chi(b)\overline{\chi(g)}}{\chi(1)}$, in any finite group.
\item Show that $\sum_{r=0}^m\frac{(-1)^r}{\binom mr}=\frac{m+1}{m+2}\bigl(1+(-1)^m\bigr)$.
\item Deduce that $c$ is the product of two $n$-cycles in exactly $\frac{2(n-1)!}{n+1}$ ways if $n$ is odd, and in none if $n$ is
even.
\end{enumerate}
\end{exercise}

\begin{exercise}[$\star\star$]\label{pb:denes}
\begin{enumerate}[label=(\alph*)]
\item Let $\tau_1,\dots,\tau_{n-1}$ be transpositions of $\{1,\dots,n\}$, $\tau_i=(a_i\,b_i)$. Show that $\tau_1\tau_2\cdots\tau_{n-1}$ is an
$n$-cycle if and only if the edges $\{a_i,b_i\}$ form a tree.
\item (Dénes, 1959.) Deduce that $(1\,2\,\dots\,n)$ can be written as a product of $n-1$ transpositions in exactly $n^{n-2}$
ways.
\end{enumerate}
\end{exercise}

\begin{figure}[ht]
\centering
\fig{ch04-tree}
\caption{A tree on $5$ vertices with its edges numbered. The product of the four transpositions, in this order, is a $5$-cycle;
each of the $4!$ orders of the edges gives a $5$-cycle, and every factorisation of a $5$-cycle into four transpositions comes
from a tree in this way (Problem~\ref{pb:denes}).}
\label{fig:tree}
\end{figure}

\section{Hints}

\begin{hint}{pb:dice}
The $24$ rotations of the cube act on numberings. Which rotation fixes a numbering?
\end{hint}

\begin{hint}{pb:patterns}
(a) On a cycle of $g$, the colouring is determined by one colour, and that colour must come back after $\ell$ steps of $s$.
Then apply Burnside's lemma to the group $G\times\mathfrak S(K)$.
\end{hint}

\begin{hint}{pb:necklacering}
(a) A word of length $n$ has a primitive period $d\mid n$. (b) Take logarithms. (c) Use (a) for $ab$, and $\gcd\cdot\operatorname{lcm}=ij$.
\end{hint}

\begin{hint}{pb:unlabelled}
Burnside, with the identity term set apart. A permutation that moves $m$ points leaves many pairs of vertices out of
its fixed pairs, and the non-fixed pairs come in cycles of length at least $2$.
\end{hint}

\begin{hint}{pb:commuting}
A $k$-tuple of commuting permutations is an action of $\Z^k$. Use the exponential formula (Tome~4, Chapter~2) on transitive
actions, and count the subgroups of index $m$ of $\Z^k$ by induction on $k$.
\end{hint}

\begin{hint}{pb:landau}
Divide the class equation by $|G|$: $1=\sum_i\frac1{|C_G(x_i)|}$, an Egyptian fraction. Order the denominators and bound
them one by one.
\end{hint}

\begin{hint}{pb:rank}
Fixed points of $g$ on $X\times X$ are pairs of fixed points. For the character, $\langle\pi,\pi\rangle=\mathbb E[F^2]$.
\end{hint}

\begin{hint}{pb:derangements}
Look at $\mathbb E[(F-1)(F-n)]$ and at the sign of $(F-1)(F-n)$ when $F\ge1$.
\end{hint}

\begin{hint}{pb:glfix}
$GL_n(\mathbb F_q)$ acts on vectors, and on pairs of vectors; count the orbits. For (b), use $q^d-1\ge q-1$ when $d\ge1$,
and Cauchy--Schwarz.
\end{hint}

\begin{hint}{pb:ksubsets}
(a) The orbits of $\mathfrak S_n$ on pairs $(A,B)$ of subsets are determined by $|A\cap B|$. (b) A virtual character of norm
$1$ and positive degree is irreducible. (c) The number of orbits of $G$ is $\langle\pi_k,1\rangle_G$.
\end{hint}

\begin{hint}{pb:samecycle}
(a) Count the permutations in which the cycle of $1$ has length $\ell$. (b) Condition on that length.
\end{hint}

\begin{hint}{pb:foata}
To go back from the word, cut it just before each record.
\end{hint}

\begin{hint}{pb:roots}
(a) Exponential formula. (b) Differentiate, and look at $(n-1)!/(n-p)!$ modulo $p$. (c) Elements of order $p$ come in groups
of $p-1$.
\end{hint}

\begin{hint}{pb:oddcycles}
The two generating functions are $\sqrt{\frac{1+x}{1-x}}$ and $\frac1{\sqrt{1-x^2}}$.
\end{hint}

\begin{hint}{pb:squares}
(a) What does squaring do to a cycle of odd length? Of even length? (b) Build the generating function cycle length by cycle
length: a factor $\exp$ for a free number of cycles, $\cosh$ for an even number.
\end{hint}

\begin{hint}{pb:poissoncycles}
(a) $\prod(C_j)_{m_j}$ counts ordered lists of distinct cycles of $\sigma$. (b) Inclusion--exclusion writes $P(C=c)$ with
factorial moments, exactly.
\end{hint}

\begin{hint}{pb:golomb}
(a) Cauchy's formula: there are $\frac{n!}{\prod j^{c_j}c_j!}$ permutations with $c_j$ cycles of length $j$. (c) Compare the sum
over $j$ with an integral, then the sum over $m$ with another integral.
\end{hint}

\begin{hint}{pb:conjugate}
(a) $P_n=\sum_\lambda z_\lambda^{-2}$, where $\frac{n!}{z_\lambda}$ is the size of the class of type $\lambda$. (b) Keep only the
types with one cycle longer than $\frac n2$.
\end{hint}

\begin{hint}{pb:commutators}
Compute $\langle N,\chi\rangle$ with Schur's lemma: $\sum_y\rho(y)M\rho(y)^{-1}$ is a scalar.
\end{hint}

\begin{hint}{pb:frobschur}
(a) $\chi(x^2)$ is the difference of the characters of $\operatorname{Sym}^2V$ and $\Lambda^2V$. (b) Compute $\langle r,\chi\rangle$.
\end{hint}

\begin{hint}{pb:ncycles}
(a) The eigenvalues of $c$, and the column orthogonality $\sum_\chi|\chi(c)|^2=|C_G(c)|$. (c) $\frac1{\binom mr}$ is a Beta
integral.
\end{hint}

\begin{hint}{pb:denes}
Multiplying by a transposition $(a\,b)$ merges two cycles if $a$ and $b$ are in different cycles, and splits one otherwise.
\end{hint}

\section{Solutions}

\begin{solution}{pb:dice}
\idea Once all six numbers are different, a rotation that fixes a numbering fixes every face, so it is the identity. The
action is free, and every orbit has $24$ elements.

\writeup (a) The rotations of the cube form a group of order $24$ (it acts transitively on the $6$ faces, and a face is fixed
by the $4$ rotations about its axis: orbit--stabiliser). It acts on the $6!=720$ numberings. A rotation that fixes a numbering
maps each face to the face with the same number, that is, to itself; the only such rotation is the identity. So all orbits
have $24$ elements, and there are $\frac{720}{24}=30$ dice.

(b) A standard numbering puts the pairs $\{1,6\}$, $\{2,5\}$, $\{3,4\}$ on the three axes ($3!$ ways) and chooses which end of
each axis gets the smaller number ($2^3$ ways): $48$ numberings, permuted by the rotations, freely as before. So there are
$\frac{48}{24}=2$ standard dice.

(c) A mirror image is $f\circ s$, where $f$ is the numbering and $s$ a reflection of the cube. If $f\circ s$ and $f$ were the
same die, then $f\circ s=f\circ g$ for some rotation $g$, so $s=g$, because $f$ is a bijection. But a reflection is not a
rotation (its determinant is $-1$). So the $30$ dice form $15$ mirror pairs, and the two standard dice are mirror images of each other.
\tools Tome~2, Chapters~1 and~3 (rotations of the cube, orbit--stabiliser).
\leads Dice are \emph{chiral}, like the molecules of Pólya's story: a molecule and its mirror image can be different
substances, as are the two forms of carvone, one smelling of spearmint and the other of caraway.
\end{solution}

\begin{solution}{pb:patterns}
\idea Renaming colours and moving points are two commuting actions, so the group is $G\times\mathfrak S(K)$, and Burnside
applies once we can count fixed colourings. Along a cycle of $g$, the condition $f(gx)=s(f(x))$ propagates one colour around,
and it closes up exactly when the colour has come back.

\writeup (a) The pair $(g,s)$ acts by $f\mapsto s\circ f\circ g^{-1}$, an action of $G\times\mathfrak S(K)$. The condition
$s\circ f\circ g^{-1}=f$ means $f(gx)=s(f(x))$ for all $x$. Let $(x,gx,\dots,g^{\ell-1}x)$ be a cycle of $g$. Then
$f(g^ix)=s^i(f(x))$, so $f$ on the cycle is determined by $c=f(x)$, and since $g^\ell x=x$, we need $s^\ell(c)=c$. Conversely
any $c$ with $s^\ell(c)=c$ gives a consistent $f$ on the cycle. Now $s^\ell(c)=c$ means that the cycle of $s$ through $c$ has
a length $j$ dividing $\ell$, and there are $\sum_{j\mid\ell}jc_j(s)$ such colours. The cycles of $g$ are independent, so the
count is the product.

(b) Here $G=\Z/n\Z$ and $K=\{0,1\}$. The rotation by $m$ has $\gcd(m,n)$ cycles of length $\ell=\frac n{\gcd(m,n)}$. For $s$
the identity, each cycle allows $2$ colours: $2^{\gcd(m,n)}$. For $s$ the exchange, $c_1(s)=0$ and $c_2(s)=1$, so a cycle allows
$2$ colours if $\ell$ is even and none otherwise. By Burnside the number of classes is
\[
\frac1{2n}\Bigl(\sum_{m=0}^{n-1}2^{\gcd(m,n)}+\sum_{\substack{0\le m<n\\ n/\gcd(m,n)\text{ even}}}2^{\gcd(m,n)}\Bigr).
\]
For $n=6$: the gcds for $m=0,\dots,5$ are $6,1,2,3,2,1$, giving $64+2+4+8+4+2=84$; the cycle length $\frac6{\gcd}$ is even for
$m=1,3,5$ (lengths $6,2,6$), giving $2+8+2=12$. So $\frac{84+12}{12}=8$.

(c) The $24$ rotations of the cube have these cycle types on the faces: the identity, $1^6$; six quarter turns about face axes,
$1^24^1$; three half turns about face axes, $1^22^2$; eight rotations about vertex axes, $3^2$; six half turns about edge axes,
$2^3$. Take $k=3$. For $s=e$, each cycle allows $3$ colours, and the sum over $G$ of $3^{c(g)}$ is
$3^6+6\cdot3^3+3\cdot3^4+8\cdot3^2+6\cdot3^3=1368$. For each of the $3$ transpositions $s$, a cycle allows $1$ colour if $\ell$ is
odd and $1+2=3$ if $\ell$ is even: the sum over $G$ is $1+6\cdot3+3\cdot9+8\cdot1+6\cdot27=216$. For each of the $2$ three-cycles
$s$, a cycle allows $3$ colours if $3\mid\ell$ and none otherwise: only the vertex rotations contribute, $8\cdot9=72$. Burnside
for the group of order $24\cdot6$:
\[
\frac{1368+3\cdot216+2\cdot72}{144}=\frac{2160}{144}=15.
\]
\tools Tome~2, Chapter~3 (Burnside's lemma, the rotations of the cube).
\leads This is de Bruijn's extension of Pólya's theorem (1959). It counts, for instance, Boolean functions up to permutation
and negation of the variables and of the value, a question of the 1950s on switching circuits.
\end{solution}

\begin{solution}{pb:necklacering}
\idea Every word of length $n$ is a primitive word of some length $d\mid n$ repeated $\frac nd$ times, and a primitive necklace
of length $d$ contains exactly $d$ words. That gives (a), and everything else is (a) in disguise: the logarithm of the product
in (b), and the factorisation $(ab)^n=a^nb^n$ in (c).

\writeup (a) A word $w$ of length $n$ has a smallest period $d$, the least $d\ge1$ such that rotating by $d$ fixes $w$. The
rotations that fix $w$ form a subgroup of $\Z/n\Z$, generated by $d$, so $d\mid n$, and $w$ is $u^{n/d}$ with $u$ a word of length
$d$ whose $d$ rotations are distinct. Conversely each such $u$ gives a word of smallest period $d$. The words $u$ of length $d$
with distinct rotations fall into primitive necklaces of exactly $d$ words. So $a^n=\sum_{d\mid n}d\,M(d,a)$, and Möbius
inversion (Tome~1, Chapter~4) gives $nM(n,a)=\sum_{d\mid n}\mu(d)a^{n/d}$.

The polynomial $P(t)=\frac1n\sum_{d\mid n}\mu(d)t^{n/d}$ has rational coefficients and takes integer values at all $t\ge1$, and also $P(0)=0$.
A polynomial with this property is an integer combination of the binomial polynomials $\binom tj$ (Newton's interpolation
formula, Tome~4, Chapter~1), and these take integer values at negative integers too. So $\sum_{d\mid n}\mu(d)a^{n/d}\equiv0\pmod n$
for every $a\in\Z$.

(b) Both sides have constant term $1$, so it is enough to compare logarithms (Tome~4, Chapter~2):
\[
\log\prod_{n\ge1}(1-x^n)^{-M(n,a)}=\sum_{n\ge1}M(n,a)\sum_{m\ge1}\frac{x^{nm}}m=\sum_{N\ge1}\frac{x^N}N\sum_{n\mid N}n\,M(n,a)
=\sum_{N\ge1}\frac{(ax)^N}N,
\]
which is $-\log(1-ax)$, by (a) with $N=nm$.

(c) Write $R(n)$ for the right-hand side. For $n\ge1$,
\[
\sum_{d\mid n}d\,R(d)=\sum_{d\mid n}\ \sum_{\operatorname{lcm}(i,j)=d}\operatorname{lcm}(i,j)\gcd(i,j)M(i,a)M(j,b)=\sum_{i\mid n,\ j\mid n}ij\,M(i,a)M(j,b),
\]
because $\operatorname{lcm}(i,j)\mid n$ if and only if $i\mid n$ and $j\mid n$, and $\operatorname{lcm}\cdot\gcd=ij$. By (a) this is
$a^nb^n=(ab)^n=\sum_{d\mid n}d\,M(d,ab)$. Möbius inversion of $d\,R(d)$ and $d\,M(d,ab)$ gives $R(n)=M(n,ab)$.

The meaning: a word over the alphabet of pairs is a pair of words, one over each alphabet. If the first has smallest period $i$
and the second $j$, the pair has smallest period $\operatorname{lcm}(i,j)$. A pair of primitive necklaces of lengths $i$ and $j$
gives $ij$ pairs of words, hence $\frac{ij}{\operatorname{lcm}(i,j)}=\gcd(i,j)$ primitive necklaces of length $\operatorname{lcm}(i,j)$.
\tools Tome~1, Chapter~4 (Möbius inversion); Tome~4, Chapters~1 and~2 (Newton series, logarithms of power series).
\leads $M(n,q)$ is also the number of monic irreducible polynomials of degree $n$ over $\mathbb F_q$ (Tome~2, Chapter~9), and then
(b) is unique factorisation in $\mathbb F_q[X]$. Identity (c) makes the sequences $(M(n,a))_n$ into a ring, the \emph{necklace
ring} of Metropolis and Rota, isomorphic to the ring of Witt vectors that controls $p$-adic arithmetic.
\end{solution}

\begin{solution}{pb:unlabelled}
\idea Burnside writes $g_n$ as an average over $\mathfrak S_n$ of $2^{c(\sigma)}$, where $c(\sigma)$ is the number of cycles of
$\sigma$ on pairs of vertices. The identity alone contributes $\frac{2^N}{n!}$. Every other permutation has far fewer cycles on
pairs, and there are not enough of them to matter.

\writeup Let $\mathfrak S_n$ act on the $N$ pairs, and $c(\sigma)$ be the number of cycles of $\sigma$ on pairs. A graph is a
$2$-colouring of the pairs, so by Burnside
\[
g_n=\frac1{n!}\sum_\sigma2^{c(\sigma)},\qquad\frac{n!\,g_n}{2^N}=1+\sum_{\sigma\ne e}2^{c(\sigma)-N}.
\]
(a) follows at once.

\emph{A bound on $c(\sigma)$.} Let $\sigma$ move exactly $m\ge2$ points. A pair fixed by $\sigma$ lies inside the $n-m$ fixed
points, or is a $2$-cycle $\{a,\sigma(a)\}$ of $\sigma$; so there are at most $f=\binom{n-m}2+\frac m2$ fixed pairs. The other
$N-f$ pairs lie in cycles of length at least $2$, so $c(\sigma)\le f+\frac{N-f}2$, and
\[
N-c(\sigma)\ge\frac{N-f}2=\frac12\Bigl(\binom n2-\binom{n-m}2-\frac m2\Bigr)=\frac{m(2n-m-2)}4.
\]

(b) and (c). There are at most $\binom nm m!\le n^m$ permutations moving $m$ points. A transposition $(a\,b)$ fixes the pair
$\{a,b\}$ and the $\binom{n-2}2$ pairs avoiding $a$ and $b$, and exchanges $\{a,x\}$ with $\{b,x\}$ for the $n-2$ other $x$: so
$N-c=n-2$ exactly, and the $\binom n2$ transpositions contribute $\binom n22^{2-n}=\frac{n(n-1)}{2^{n-1}}$. For $m=3$, the $\sigma$ contribute at most
$n^32^{-3(2n-5)/4}=O(n^32^{-3n/2})$. For $4\le m\le\frac n2$, $2n-m-2\ge\frac{3n-4}2$, so they contribute at most
$\bigl(n2^{-(3n-4)/8}\bigr)^m$, a geometric series in $m$ whose sum over $m\ge4$ is $O(n^42^{-3n/2})$. For $m>\frac n2$, $2n-m-2\ge n-2$, so $N-c\ge\frac{n(n-2)}8$,
and the at most $n!$ such $\sigma$ contribute at most $n^n2^{-n(n-2)/8}$, which is far smaller. Altogether
\[
\frac{n!\,g_n}{2^N}=1+\frac{n(n-1)}{2^{n-1}}+O\bigl(n^42^{-3n/2}\bigr)\longrightarrow1.
\]
\tools Tome~2, Chapter~3 (Burnside's lemma); Tome~4, Chapter~11 (estimating sums).
\leads The meaning of $g_n\sim\frac{2^N}{n!}$ is that almost every graph has a trivial automorphism group (Erdős and Rényi,
1963). The full asymptotic expansion was found by Oberschelp (1967). Numerical evaluation shows how fast the ratio converges: at $n=10$ it is
already $1.24$, and the two-term formula predicts $1.18$.
\end{solution}

\begin{solution}{pb:commuting}
\idea Commuting permutations $\sigma_1,\dots,\sigma_k$ are the same thing as an action of $\Z^k$ on $\{1,\dots,n\}$, the $i$-th
generator acting by $\sigma_i$. An action splits into transitive pieces, and the exponential formula turns ``a set of
transitive pieces'' into an exponential. A transitive action of an abelian group on $m$ points is a subgroup of index $m$ plus a
labelling. So everything reduces to counting subgroups of $\Z^k$ of index $m$.

\writeup Let $\Gamma=\Z^k$. Commuting $k$-tuples of permutations of a set $S$ correspond to actions of $\Gamma$ on $S$. An action
is the disjoint union of its restrictions to its orbits, which are transitive actions. Let $t_m$ be the number of transitive
actions of $\Gamma$ on $\{1,\dots,m\}$. By the exponential formula (Tome~4, Chapter~2),
\[
\sum_{n\ge0}h_k(n)\frac{x^n}{n!}=\exp\Bigl(\sum_{m\ge1}t_m\frac{x^m}{m!}\Bigr).
\]
\emph{Transitive actions.} Given a transitive action on $\{1,\dots,m\}$, let $H$ be the stabiliser of the point $1$: a subgroup
of index $m$ (orbit--stabiliser), and $\gamma H\mapsto\gamma\cdot1$ is a bijection $\Gamma/H\to\{1,\dots,m\}$ sending $H$ to $1$.
Conversely a subgroup $H$ of index $m$ and a bijection $\beta:\Gamma/H\to\{1,\dots,m\}$ with $\beta(H)=1$ define a transitive action
$\gamma\cdot\beta(\delta H)=\beta(\gamma\delta H)$, whose stabiliser of $1$ is $H$. These are inverse constructions, and there are
$(m-1)!$ bijections $\beta$. So $t_m=(m-1)!\,a_m(\Gamma)$, with $a_m$ the number of subgroups of index $m$, and
\[
\sum_nh_k(n)\frac{x^n}{n!}=\exp\Bigl(\sum_{m\ge1}a_m(\Z^k)\frac{x^m}m\Bigr).
\]
\emph{Subgroups of $\Z^k$ of index $m$.} We show $a_m(\Z^k)=\sum_{e\mid m}e\,a_e(\Z^{k-1})$. Let $H\le\Z^k$ have index $m$, write
$\Z^k=\Z^{k-1}\times\Z$, let $f\Z$ be the image of $H$ under the last coordinate ($f\ge1$ since the index is finite) and
$H'=H\cap(\Z^{k-1}\times\{0\})$. Then $H=H'\oplus\Z v$ for any $v=(w,f)\in H$, the index of $H'$ in $\Z^{k-1}$ is $e=\frac mf$, and $w$ is
determined modulo $H'$ by $H$. Conversely, for $f\mid m$, $H'$ of index $\frac mf$ in $\Z^{k-1}$ and any $w$, the group
$H'\oplus\Z(w,f)$ has last-coordinate image $f\Z$, meets $\Z^{k-1}\times\{0\}$ in $H'$ (an element $h'+t(w,f)$ has last coordinate
$tf$), and has index $\frac mf\cdot f=m$; and two choices $w,w'$ give the same group if and only if $w-w'\in H'$, which leaves
$e=\frac mf$ choices. Hence the recursion. With $a_m(\Z)=1$: $a_m(\Z^2)=\sigma(m)$ and $a_m(\Z^3)=\sum_{e\mid m}e\,\sigma(e)$.

\emph{The products.} For $k=2$, writing $m=de$,
\[
\sum_m\sigma(m)\frac{x^m}m=\sum_{d,e\ge1}d\,\frac{x^{de}}{de}=\sum_{e\ge1}-\log(1-x^e),
\]
so the exponential is $\prod_e(1-x^e)^{-1}$. For $k=3$, writing $m=de$,
\[
\sum_m\Bigl(\sum_{d\mid m}d\,\sigma(d)\Bigr)\frac{x^m}m=\sum_{d,e}\sigma(d)\frac{x^{de}}e=\sum_d-\sigma(d)\log(1-x^d),
\]
so the exponential is $\prod_d(1-x^d)^{-\sigma(d)}$. The first product is $\sum p(n)x^n$ (Tome~4, Chapter~2), so $h_2(n)=n!p(n)$.
\tools Tome~2, Chapters~3 and~7 (orbit--stabiliser, subgroups of $\Z^k$); Tome~4, Chapter~2 (exponential formula, partitions); Tome~1, Chapter~4 (divisor sums).
\leads The same argument gives, for every $k$,
\[
\sum_nh_k(n)\frac{x^n}{n!}=\exp\Bigl(\sum_ma_m(\Z^k)\frac{x^m}m\Bigr),\qquad\sum_m\frac{a_m(\Z^k)}{m^s}=\zeta(s)\zeta(s-1)\cdots\zeta(s-k+1).
\] Bryan and Fulman (1998) found these products as ``orbifold Euler
characteristics'' in string theory. $h_2(n)=n!p(n)$ also follows from Tome~2, Exercise~3.8: the number of
commuting pairs is $|G|$ times the number of classes, and $\mathfrak S_n$ has $p(n)$ classes.
\end{solution}

\begin{solution}{pb:landau}
\idea The class equation divided by $|G|$ is an Egyptian fraction: $1=\sum_{i=1}^k\frac1{c_i}$, $c_i$ the order of a
centraliser. The largest $c_i$ is $|G|$ itself (the centraliser of $e$). With $k$ terms one cannot reach far: the smallest
denominator is at most $k$, and each next one is bounded by the previous ones.

\writeup (a) Let $x_1,\dots,x_k$ represent the classes. The class of $x_i$ has $\frac{|G|}{c_i}$ elements with
$c_i=|C_G(x_i)|$ (orbit--stabiliser for conjugation), so $1=\sum_i\frac1{c_i}$, and $c_i=|G|$ for the class of $e$. Order
$c_1\le c_2\le\dots\le c_k$; then $c_k=|G|$. Let $r_i=1-\sum_{j\le i}\frac1{c_j}$. For $0\le i<k$, $r_i=\sum_{j>i}\frac1{c_j}\le\frac{k-i}{c_{i+1}}$,
and $r_i>0$ is a fraction with denominator dividing $c_1\cdots c_i$, so $r_i\ge\frac1{c_1\cdots c_i}$. Hence
$c_{i+1}\le(k-i)c_1\cdots c_i\le k\,c_1\cdots c_i$. By induction $c_i\le k^{2^{i-1}}$: true for $i=1$ ($c_1\le k$), and then
$c_{i+1}\le k\cdot k^{1+2+\dots+2^{i-1}}=k^{2^i}$. So $|G|=c_k\le k^{2^{k-1}}$. There are finitely many groups of each order,
hence finitely many with $k$ classes.

(b) $k=1$: $G=\{e\}$. $k=2$: $1=\frac1{c_1}+\frac1{c_2}$ with $c_1\le c_2$ forces $c_1=c_2=2$, so $|G|=2$ and $G=\Z/2\Z$. $k=3$: the
solutions of $1=\frac1{c_1}+\frac1{c_2}+\frac1{c_3}$, $c_1\le c_2\le c_3$, are $(3,3,3)$, $(2,4,4)$, $(2,3,6)$, so $|G|\in\{3,4,6\}$.
Order $3$: $\Z/3\Z$, which has $3$ classes. Order $4$: abelian, so $4$ classes, impossible. Order $6$: $\Z/6\Z$ has $6$ classes,
and $\mathfrak S_3$ has $3$. So the answer is $\{e\}$, $\Z/2\Z$, $\Z/3\Z$ and $\mathfrak S_3$.
\tools Tome~2, Chapter~3 (class equation); Tome~1, Chapter~1 (fractions).
\leads This is Landau's theorem (1903). The sharp bound for the largest denominator of an Egyptian fraction with $k$ terms is
$s_{k-1}-1$, where $s_0=2$, $s_1=3$, $7$, $43$, \dots{} is Sylvester's sequence of Problem~\ref{pb:sylvester} (Curtiss, 1922); it gives
$|G|\le s_{k-1}-1$, which is doubly exponential too. The groups with few classes have been classified, with the help of computers, up
to about a dozen classes.
\end{solution}

\begin{solution}{pb:rank}
\idea A pair $(x,y)$ is fixed by $g$ exactly when both points are, so $F^2$ counts the fixed points of $g$ on $X\times X$, and
Burnside turns its mean into a number of orbits. The diagonal is one orbit; the rest is at least one more.

\writeup (a) By Burnside, $\mathbb E[F]$ is the number of orbits of $G$ on $X$, which is $1$. The fixed points of $g$ on $X\times X$
are the pairs of fixed points, $F(g)^2$ of them, so $\mathbb E[F^2]=r$. The diagonal $\{(x,x)\}$ is an orbit, because $G$ is
transitive, and its complement is non-empty since $n\ge2$; so $r\ge2$.

(b) $r=2$ if and only if the off-diagonal pairs form a single orbit, which is $2$-transitivity. $\operatorname{Var}F=r-1$, so
this is also $\operatorname{Var}F=1$. Finally $\pi$ is real-valued, so $\langle\pi,\pi\rangle=\frac1{|G|}\sum\pi(g)^2=r$ and
$\langle\pi,1\rangle=\mathbb E[F]=1$. With $\psi=\pi-1$, $\langle\psi,1\rangle=0$ and $\langle\psi,\psi\rangle=r-2\langle\psi,1\rangle-1=r-1$.
A character $\psi=\sum m_i\chi_i$ has $\langle\psi,\psi\rangle=\sum m_i^2$, so it is irreducible if and only if
$\langle\psi,\psi\rangle=1$ (Tome~2, Chapter~11), that is, $r=2$.

(c) A rotation or reflection of the $n$-gon preserves the distance between two vertices along the polygon, which takes the
values $0,1,\dots,\lfloor\frac n2\rfloor$. Conversely, if $(x,y)$ and $(x',y')$ are at the same distance, a rotation takes $x$ to $x'$
and then $y$ to $y'$ or to the mirror image of $y'$ in the axis through $x'$, and the reflection in that axis repairs it. So
the orbits on pairs are the distance classes, and $r=\lfloor\frac n2\rfloor+1$. For $n=3$ this is $2$: $D_3=\mathfrak S_3$ is
$2$-transitive on the triangle.
\tools Tome~2, Chapters~3 and~11 (Burnside, inner products of characters); Tome~3, Chapter~2 (variance).
\leads Problem~\ref{pb:fixmoments} is the case $G=\mathfrak S_n$, where all moments are Bell numbers. The higher moments
$\mathbb E[F^j]$ count orbits on $X^j$, and a group is $j$-transitive exactly when they agree with those of $\mathfrak S_n$ up to
order $j$.
\end{solution}

\begin{solution}{pb:derangements}
\idea We know $\mathbb E[F]=1$ and $\mathbb E[F^2]=r\ge2$. Find a quadratic function of $F$ that is $\le0$ whenever $F\ge1$ and
positive at $F=0$: $(F-1)(F-n)$ is the natural choice, since $1\le F\le n$ when $F\ge1$. Its mean is $r-1\ge1$, and only the
derangements can produce it.

\writeup (a) Let $D$ be the set of $g$ with $F(g)=0$. For $g\notin D$, $1\le F(g)\le n$, so $(F-1)(F-n)\le0$. For $g\in D$ it is $n$.
So
\[
n\,P(D)\ge\mathbb E[(F-1)(F-n)]=\mathbb E[F^2]-(n+1)\mathbb E[F]+n=r-1\ge1,
\]
by Problem~\ref{pb:rank}. Hence $P(D)\ge\frac1n$.

(b) Equality needs $r=2$ and $(F-1)(F-n)=0$ for every $g\notin D$, that is $F(g)\in\{0,1,n\}$ for all $g$. Since the action is
faithful, $F(g)=n$ only for $g=e$. So equality holds if and only if $G$ is $2$-transitive and every $g\ne e$ fixes at most one
point. If so, the stabiliser of two distinct points is trivial, so the $g$ with $gx=x'$, $gy=y'$ is unique: $G$ is sharply
$2$-transitive. Conversely, sharply $2$-transitive means $2$-transitive with trivial two-point stabilisers, so every $g\ne e$
fixes at most one point.

For $t\mapsto at+b$ on $\mathbb F_p$: given $x\ne y$ and $x'\ne y'$, the system $ax+b=x'$, $ay+b=y'$ has the unique solution
$a=\frac{x'-y'}{x-y}\ne0$, $b=x'-ax$. So the group, of order $p(p-1)$, is sharply $2$-transitive, and its derangements are the
$p-1$ translations $t\mapsto t+b$, $b\ne0$: a proportion $\frac{p-1}{p(p-1)}=\frac1p$.

(c) Let $H<G$ be proper and suppose $G=\bigcup_{g\in G}gHg^{-1}$. Then $G$ acts transitively on the
$n=[G:H]\ge2$ cosets $gH$, and $x$ fixes $gH$ if and only if $g^{-1}xg\in H$, that is $x\in gHg^{-1}$. So every element would fix a
coset. The action may not be faithful, but (a) applies to the faithful action of $G/K$, $K$ the kernel, and gives an element
with no fixed point. Contradiction.
\tools Tome~2, Chapter~3 (actions on cosets); Tome~3, Chapter~2 (expectations).
\leads Jordan's theorem (1872) has a famous arithmetic consequence: an irreducible polynomial of degree at least $2$ in
$\Z[X]$ has no root modulo $p$ for infinitely many primes $p$, a set of positive density by Chebotarev's theorem (see Serre's
article below). Fulman and Guralnick proved a conjecture of Boston and Shalev: for simple groups the proportion of
derangements is bounded below by an absolute constant.
\end{solution}

\begin{solution}{pb:glfix}
\idea Fixed vectors of $A$ form the subspace $\ker(A-I)$, with $q^d$ elements. So $q^d$ is the number of fixed points of $A$ on
$V=\mathbb F_q^n$, and $q^{2d}$ the number on $V\times V$. Count the orbits of $GL_n$ on vectors and on pairs of vectors.

\writeup (a) $GL_n(\mathbb F_q)$ acts on $V$ with two orbits, $\{0\}$ and $V\setminus\{0\}$ (any non-zero vector can be completed
to a basis, and an invertible matrix maps one basis to another). By Burnside $\mathbb E[q^d]=2$. On $V\times V$, the orbit of
$(v,w)$ is determined by the linear relations between $v$ and $w$: the orbits are $\{(0,0)\}$; $\{(0,w):w\ne0\}$; $\{(v,0):v\ne0\}$;
$\{(v,\lambda v):v\ne0\}$ for each $\lambda\in\mathbb F_q^*$; and the pairs of independent vectors, one orbit since $n\ge2$ (complete
$(v,w)$ to a basis). That is $3+(q-1)+1=q+3$ orbits, so $\mathbb E[q^{2d}]=q+3$.

(b) Let $Y=q^d-1\ge0$. Then $\mathbb E[Y]=1$ and $\mathbb E[Y^2]=(q+3)-2\cdot2+1=q$. When $d\ge1$, $Y\ge q-1$, so
$1=\mathbb E[Y]\ge(q-1)P(d\ge1)$. And by Cauchy--Schwarz, $1=\mathbb E[Y\mathbf 1_{d\ge1}]^2\le\mathbb E[Y^2]P(d\ge1)=q\,P(d\ge1)$.
\tools Tome~0, Chapter~7 (linear algebra); Tome~2, Chapter~3 (Burnside); Tome~3, Chapter~2 (Cauchy--Schwarz).
\leads For $q=2$ the bounds are $\frac12$ and $1$, and one finds numerically $0.714$ for $n=3$. The exact law of $d$ comes from the
\emph{cycle index} of $GL_n(\mathbb F_q)$, a $q$-analogue of the cycle index of $\mathfrak S_n$ (see Fulman's survey). As
$n\to\infty$ with $q$ fixed, the probability of the eigenvalue $1$ converges, and the limit tends to $0$ like $\frac1q$.
\end{solution}

\begin{solution}{pb:ksubsets}
\idea Inner products of permutation characters count orbits on pairs, and for pairs of subsets the only invariant is the size of
the intersection. So the matrix of inner products is $\min(k,l)+1$, the matrix of the ``staircase'' $\pi_k=\chi_0+\dots+\chi_k$
with orthonormal $\chi_j$. The Livingstone--Wagner theorem then falls out by restriction to $G$.

\writeup (a) $\pi_k\pi_l(\sigma)$ is the number of pairs $(A,B)$, $|A|=k$, $|B|=l$, fixed by $\sigma$, so
$\langle\pi_k,\pi_l\rangle$ is the number of orbits of $\mathfrak S_n$ on such pairs ($\pi_k$ is real-valued). If $|A\cap B|=|A'\cap B'|=i$,
the four regions $A\cap B$, $A\setminus B$, $B\setminus A$ and the complement of $A\cup B$ have the same sizes for both pairs, so a
permutation maps one pair to the other. Hence the orbits correspond to the possible values of $i$, which are
$\max(0,k+l-n),\dots,\min(k,l)$; for $k,l\le\frac n2$ that is $0,\dots,\min(k,l)$, and there are $\min(k,l)+1$.

(b) With $\pi_{-1}=0$, for $0\le k\le l\le\frac n2$:
\[
\langle\chi_k,\chi_l\rangle=\langle\pi_k,\pi_l\rangle-\langle\pi_{k-1},\pi_l\rangle-\langle\pi_k,\pi_{l-1}\rangle+\langle\pi_{k-1},\pi_{l-1}\rangle,
\]
which is $(k+1)-k-(k+1)+k=0$ if $k<l$ and $(k+1)-k-k+k=1$ if $k=l$ (for $k=0$ drop the terms with $\pi_{-1}$; the result is the
same). Each $\chi_k=\sum_im_i\psi_i$ is an integer combination of irreducible characters $\psi_i$ with $\sum m_i^2=1$, so
$\chi_k=\pm\psi$ for one irreducible $\psi$; and $\chi_k(e)=\binom nk-\binom n{k-1}>0$ for $k\le\frac n2$, so $\chi_k=\psi$. They are
distinct because they are orthogonal.

(c) The number of orbits of $G$ on $k$-subsets is $\frac1{|G|}\sum_{g\in G}\pi_k(g)=\langle\pi_k|_G,1\rangle_G$ by Burnside. By (b),
$\pi_k=\chi_0+\chi_1+\dots+\chi_k$, so this number is $\sum_{j\le k}\langle\chi_j|_G,1\rangle_G$. Each $\chi_j$ is the character of a
representation of $\mathfrak S_n$, hence $\chi_j|_G$ of a representation of $G$, and $\langle\chi_j|_G,1\rangle_G$ is the
dimension of its fixed vectors, at least $0$. So the sum is non-decreasing in $k$.
\tools Tome~2, Chapters~3 and~11 (Burnside, orthogonality of characters); Tome~0, Chapter~1 (binomial coefficients).
\leads The $\chi_k$ are the characters of the partitions $(n-k,k)$ in Young's theory. For $n=4$ they have degrees $1,3,2$; for
$n=5$, $1,4,5$. Livingstone and Wagner (1965) proved (c) to show that a group that is transitive on $k$-subsets is also transitive
on $(k-1)$-subsets.
\end{solution}

\begin{solution}{pb:samecycle}
\idea Building the cycle of $1$ step by step: it has length $\ell$ for $(n-1)(n-2)\cdots(n-\ell+1)\cdot(n-\ell)!=(n-1)!$
permutations, whatever $\ell$ is. Then, given the length, the other points of the cycle are a uniform random subset.

\writeup (a) A permutation in which the cycle of $1$ has length $\ell$ is given by the sequence $\sigma(1),\sigma^2(1),\dots,\sigma^{\ell-1}(1)$
of distinct elements different from $1$, in $(n-1)\cdots(n-\ell+1)$ ways, and a permutation of the remaining $n-\ell$ points, in
$(n-\ell)!$ ways. That is $(n-1)!$, so the length is uniform.

(b) Given that the cycle of $1$ has length $\ell$, by symmetry its other $\ell-1$ points are a uniform $(\ell-1)$-subset of
$\{2,\dots,n\}$, which contains $\{2,\dots,k\}$ with probability $\binom{n-k}{\ell-k}/\binom{n-1}{\ell-1}=\binom{\ell-1}{k-1}/\binom{n-1}{k-1}$.
Averaging over $\ell$ and using $\sum_{\ell=k}^n\binom{\ell-1}{k-1}=\binom nk$:
\[
P=\frac1n\sum_{\ell=1}^n\frac{\binom{\ell-1}{k-1}}{\binom{n-1}{k-1}}=\frac{\binom nk}{n\binom{n-1}{k-1}}=\frac1k.
\]
\tools Tome~3, Chapters~1 and~2 (conditioning); Tome~0, Chapter~1 (the hockey-stick identity).
\leads The answer does not depend on $n$, a sign of the \emph{Chinese restaurant process}: customers $1,2,\dots$ arrive, and each
joins an existing cycle or opens a new one, which builds a random permutation one point at a time (Tome~3, Theorem~2.12).
\end{solution}

\begin{solution}{pb:foata}
\idea In the word, the first letters of the cycles are exactly the records, because each cycle starts with its maximum and
the maxima increase. So the cycles can be read back by cutting before each record.

\writeup (a) In the standard form, the first letter of each cycle is its largest element, and these first letters increase from
cycle to cycle. Each later letter of a cycle is smaller than the cycle's first letter, which is the largest letter seen so far.
So the records of the word are exactly the first letters of the cycles. Conversely, given a word, cut it before each record:
the pieces start with their largest element and their first letters increase, so reading the pieces as cycles gives a
permutation whose standard form is the word. The two maps are inverse to each other, and the number of cycles is the number of
records.

(b) A bijection maps the uniform law to the uniform law, so the two counts have the same law. The number of permutations with
$k$ cycles is $\stirlingi nk$ by definition, and the bijection transports it to the words.
\tools Tome~2, Chapter~4 (cycle decomposition); Tome~3, Chapter~2 (records, Feller's coupling).
\leads This is Foata's \emph{fundamental transformation} (Foata and Schützenberger, 1970). It explains why Tome~3 found the same
law, a sum of independent Bernoulli$(\frac1j)$, for the records (Exercise~2.1) and for the cycles
(Theorem~2.12).
\end{solution}

\begin{solution}{pb:roots}
\idea $\sigma^k=e$ means that every cycle length divides $k$, a restriction on the cycle lengths, which the exponential formula
handles. Differentiating gives a recursion, and modulo $p$ the coefficient $\frac{(n-1)!}{(n-p)!}$, a product of $p-1$
consecutive integers, is $0$ or $-1$.

\writeup (a) A permutation is a set of cycles, and there are $(d-1)!$ cycles on a $d$-set. By the exponential formula
(Tome~4, Chapter~2), permutations whose cycle lengths lie in a set $D$ have exponential generating function
$\exp\bigl(\sum_{d\in D}(d-1)!\frac{x^d}{d!}\bigr)=\exp\bigl(\sum_{d\in D}\frac{x^d}d\bigr)$. Take $D$ the divisors of $k$.

(b) For $k=p$, $A(x)=\exp(x+\frac{x^p}p)$ and $A'=(1+x^{p-1})A$. Comparing coefficients of $\frac{x^{n-1}}{(n-1)!}$:
$a_n=a_{n-1}+\frac{(n-1)!}{(n-p)!}a_{n-p}$. The product $(n-p+1)\cdots(n-1)$ runs over $p-1$ consecutive integers, which cover all
residues modulo $p$ except that of $n$. If $p\nmid n$, one of them is divisible by $p$ and the product is $\equiv0$; if $p\mid n$,
they are the non-zero residues and the product is $\equiv(p-1)!\equiv-1$ (Wilson). So $a_n\equiv a_{n-1}$ if $p\nmid n$ and
$a_n\equiv a_{n-1}-a_{n-p}$ if $p\mid n$. Since $a_0=\dots=a_{p-1}=1$, $a_p\equiv1-1=0$, and by induction all later terms vanish
modulo $p$: if $p\nmid n$, $a_n\equiv a_{n-1}\equiv0$; if $p\mid n$ and $n\ge2p$, $a_n\equiv0-0$.

(c) The solutions of $\sigma^p=e$ other than $e$ have order $p$. Each generates a cyclic group of order $p$, which contains
$p-1$ elements of order $p$, and two such groups are equal or meet in $\{e\}$. So $a_n-1$ is a multiple of $p-1$. With
$p\mid a_n$ and $\gcd(p,p-1)=1$, $a_n\equiv p\pmod{p(p-1)}$, since $p$ satisfies both congruences.
\tools Tome~4, Chapter~2 (exponential formula); Tome~1, Chapter~3 (Wilson); Tome~2, Chapter~1 (cyclic subgroups).
\leads Frobenius (1895) proved that in every finite group the number of solutions of $x^m=e$ is a multiple of
$\gcd(m,|G|)$; (b) is the case $\mathfrak S_n$, $m=p$. A famous conjecture of Frobenius, proved in 1991 with the
classification of finite simple groups, says that if there are exactly $m$ solutions, they form a subgroup.
\end{solution}

\begin{solution}{pb:oddcycles}
\idea Restricting cycle lengths to odd or even numbers gives $\exp\sum_{j\text{ odd}}\frac{x^j}j=\sqrt{\frac{1+x}{1-x}}$ and
$\exp\sum_{j\text{ even}}\frac{x^j}j=\frac1{\sqrt{1-x^2}}$, and $\sqrt{\frac{1+x}{1-x}}=\frac{1+x}{\sqrt{1-x^2}}$. The even parts agree.

\writeup (a) By the exponential formula as in Problem~\ref{pb:roots},
\[
\sum o_n\frac{x^n}{n!}=\exp\Bigl(\sum_{j\text{ odd}}\frac{x^j}j\Bigr)=\exp\Bigl(\frac12\log\frac{1+x}{1-x}\Bigr)=\frac{1+x}{\sqrt{1-x^2}},
\]
and
\[
\sum e_n\frac{x^n}{n!}=\exp\Bigl(\sum_{j\text{ even}}\frac{x^j}j\Bigr)=\exp\Bigl(-\frac12\log(1-x^2)\Bigr)=\frac1{\sqrt{1-x^2}}.
\]
Now $\frac1{\sqrt{1-x^2}}=\sum_m\binom{2m}m\frac{x^{2m}}{4^m}$ (Tome~3, Exercise~3.7). So for $n=2m$,
\[
o_n=e_n=(2m)!\binom{2m}m4^{-m}=\Bigl(\frac{(2m)!}{2^mm!}\Bigr)^2=\bigl(1\cdot3\cdots(2m-1)\bigr)^2,
\]
and the coefficient of $x^{2m+1}$ in $\frac{1+x}{\sqrt{1-x^2}}$ is $\binom{2m}m4^{-m}$ too, so $o_{2m+1}=(2m+1)!\binom{2m}m4^{-m}=(2m+1)o_{2m}$.

(b) The order of $\sigma$ is the lcm of its cycle lengths, odd exactly when all cycles are odd. The probability is
$\frac{o_n}{n!}=\binom{2m}m4^{-m}$ with $m=\lfloor\frac n2\rfloor$, and $\binom{2m}m4^{-m}\sim\frac1{\sqrt{\pi m}}$ (Stirling, Tome~4,
Chapter~11), which is $\sim\sqrt{\frac2{\pi n}}$.
\tools Tome~4, Chapters~2 and~11 (exponential formula, Stirling); Tome~3, Chapter~3 (the binomial series).
\leads For even $n$ the equality $o_n=e_n$ cries out for a bijection between permutations with all cycles odd and permutations
with all cycles even. Several are known, none obvious; finding one is a good challenge.
\end{solution}

\begin{solution}{pb:squares}
\idea Squaring keeps an odd cycle as one cycle of the same length and breaks an even cycle into two equal halves. So even
cycles of a square come in pairs. That is a condition on the cycle type, cycle length by cycle length, and the exponential
formula still works if, for even lengths, we allow only even numbers of cycles.

\writeup (a) Let $\tau$ have a cycle $(a_1\,a_2\,\dots\,a_\ell)$. If $\ell$ is odd, $\tau^2$ maps $a_i\mapsto a_{i+2}$, and since $2$ is
invertible modulo $\ell$ this is a single $\ell$-cycle. If $\ell=2j$, $\tau^2$ splits it into $(a_1\,a_3\,\dots)$ and
$(a_2\,a_4\,\dots)$, two $j$-cycles. So in $\sigma=\tau^2$, the cycles of an even length $j$ come only from $2j$-cycles of $\tau$,
two at a time: their number is even. Conversely, suppose the condition holds. An odd cycle $c$ of length $\ell$ is a square:
$c=(c^{(\ell+1)/2})^2$. Two cycles $(a_1\,\dots\,a_j)$ and $(b_1\,\dots\,b_j)$ of the same length are the square of the
$2j$-cycle $(a_1\,b_1\,a_2\,b_2\,\dots\,a_j\,b_j)$. Group the even cycles of $\sigma$ in pairs of equal length, take these roots,
and multiply them (they have disjoint supports, so they commute).

(b) A permutation is determined by its set of cycles, and by the exponential formula with cycle lengths treated separately,
the generating function of permutations with $m_\ell$ cycles of length $\ell$ allowed only for $m_\ell$ in a set $S_\ell$ is
$\prod_\ell\sum_{m\in S_\ell}\frac1{m!}\bigl(\frac{x^\ell}\ell\bigr)^m$. (The coefficient of $\frac{x^n}{n!}$ counts permutations with
those cycle types, by Cauchy's formula $\frac{n!}{\prod\ell^{m_\ell}m_\ell!}$.) For odd $\ell$ every $m$ is allowed, giving
$\exp\frac{x^\ell}\ell$; for even $\ell$ only even $m$, giving $\cosh\frac{x^\ell}\ell$. The odd factors multiply to $\sqrt{\frac{1+x}{1-x}}$
as in Problem~\ref{pb:oddcycles}.

(c) $\sqrt{\frac{1+x}{1-x}}=\frac{1+x}{\sqrt{1-x^2}}$, and the other factors are even functions of $x$. So the generating
function is $(1+x)E(x)$ with $E$ even, and the coefficients of $x^{2m}$ and $x^{2m+1}$ are equal: $\frac{s_{2m+1}}{(2m+1)!}=\frac{s_{2m}}{(2m)!}$.
And a permutation with all cycles odd is a square by (a), so $s_n\ge o_n$.
\tools Tome~2, Chapter~4 (cycles, powers); Tome~4, Chapter~2 (exponential formula).
\leads Blum (1974) proved that the probability is $\sim\sqrt{\frac2{\pi n}}\prod_{m\ge1}\cosh\frac1{2m}=1.2218\sqrt{\frac2{\pi n}}$. The
a direct computation verifies this up to $n=800$, where the ratio is $1.002$. Pouyanne (2002) treated $k$-th roots.
\end{solution}

\begin{solution}{pb:poissoncycles}
\idea A falling factorial $(C_j)_m$ counts ordered lists of $m$ distinct $j$-cycles of $\sigma$. So the mixed factorial moment
counts ordered lists of disjoint cycles that are all present in $\sigma$, and each list is present with probability
$\frac{(n-M)!}{n!}$, $M$ the number of points it uses. Since the factorial moments are exact, inclusion--exclusion gives the law
exactly, as a truncated Poisson series.

\writeup (a) Let $M=\sum jm_j$. The product $\prod_j(C_j)_{m_j}$ is the number of ways to choose, for each $j$, an ordered list of
$m_j$ distinct $j$-cycles of $\sigma$; all these cycles are disjoint. If $M>n$ there are none. Otherwise, the expectation is the
number of such lists of disjoint cycles in $\{1,\dots,n\}$ times the probability that a given list consists of cycles of $\sigma$.
A list is given by an ordered sequence of $M$ distinct points, cut into consecutive blocks of the prescribed lengths, each block
read as a cycle; a $j$-cycle can be written in $j$ ways, so there are $\frac{n!}{(n-M)!\prod_jj^{m_j}}$ lists. A given list of
disjoint cycles is in $\sigma$ for $(n-M)!$ permutations. The product is $\prod_jj^{-m_j}$.

(b) For a vector $C=(C_1,\dots,C_k)$ of bounded non-negative integers and $c\in\N^k$, inclusion--exclusion gives
\[
P(C=c)=\sum_{m\ge c}\ \prod_j(-1)^{m_j-c_j}\binom{m_j}{c_j}\ \mathbb E\Bigl[\prod_j\frac{(C_j)_{m_j}}{m_j!}\Bigr],
\]
a finite sum: expand $\mathbf 1_{C_j=c_j}=\binom{C_j}{c_j}(1-1)^{C_j-c_j}$ and use $\binom{C_j}{c_j}\binom{C_j-c_j}{i}=\binom{c_j+i}{c_j}\binom{C_j}{c_j+i}$.
By (a), the terms are those of
\[
\prod_j\sum_{m_j\ge c_j}(-1)^{m_j-c_j}\binom{m_j}{c_j}\frac{j^{-m_j}}{m_j!}=\prod_je^{-1/j}\frac{j^{-c_j}}{c_j!},
\]
restricted to $\sum jm_j\le n$. The full series converges absolutely (its absolute values sum to $\prod_j\frac{e^{1/j}j^{-c_j}}{c_j!}$), so
the omitted tail tends to $0$ as $n\to\infty$, and $P(C=c)$ tends to the product of the Poisson probabilities.

(c) By (b), $P(C_1=\dots=C_{m-1}=0)\to\prod_{j<m}e^{-1/j}=e^{-H_{m-1}}$. As $H_{m-1}=\log m+\gamma+O(\frac1m)$, this is
$\sim\frac{e^{-\gamma}}m$. For integers, the density of those with no prime factor below $m$ is $\prod_{p<m}(1-\frac1p)\sim\frac{e^{-\gamma}}{\log m}$
(Mertens, Tome~5). The same constant $e^{-\gamma}$ appears, with $m$ in place of $\log m$: cycles of length $j$ play the role of primes
of size about $e^j$, and $\frac1j$ the role of $\frac1p$.
\tools Tome~3, Chapters~2 and~6 (factorial moments, convergence in law); Tome~5, Chapter~1 (Mertens).
\leads The convergence is extremely fast: Arratia and Tavaré (1992) showed that the total variation distance is
superexponentially small in $\frac nk$. The analogy between cycles and prime factors is the theme of the book by Arratia, Barbour
and Tavaré, and it returns in Chapter~13 with Dickman's function.
\end{solution}

\begin{solution}{pb:golomb}
\idea The obstacle is that the cycle counts of a permutation of fixed size are dependent, since $\sum jC_j=n$. Randomise the size
geometrically, and the dependence disappears: the counts become exactly independent Poisson variables. The longest cycle is then
easy to describe, and the average over $n$ is a generating function in $x$ that can be studied as $x\to1$.

\writeup (a) By Cauchy's formula there are $\frac{n!}{\prod j^{c_j}c_j!}$ permutations of type $c$, $\sum jc_j=n$. So
\[
P(C=c)=(1-x)x^n\prod_j\frac1{j^{c_j}c_j!}=(1-x)\prod_j\frac{(x^j/j)^{c_j}}{c_j!}=\prod_je^{-x^j/j}\frac{(x^j/j)^{c_j}}{c_j!},
\]
because $1-x=\exp(-\sum_jx^j/j)$. A joint law that factorises in this way is that of independent Poisson$(\frac{x^j}j)$
variables (summing over all $c$ gives $1$, so every $c$ with finitely many non-zero terms is accounted for).

(b) The longest cycle $L$ is at most $m$ if and only if $C_j=0$ for all $j>m$, so by (a)
$P(L\le m)=\prod_{j>m}e^{-x^j/j}=\exp(-\sum_{j>m}\frac{x^j}j)$. Then $\mathbb E[L]=\sum_{m\ge0}P(L>m)$, and conditioning on $N$
gives the left-hand side $\sum_n(1-x)x^n\mathbb E[L_n]$.

(c) Let $f(v)=\frac{e^{-v/t}}v$, decreasing on $]0,\infty[$, with $\int_m^\infty f=E_1(\frac mt)$. For $m\ge1$, comparing sum and
integral, $E_1(\frac mt)-f(m)\le\sum_{j>m}f(j)\le E_1(\frac mt)$. The function $u\mapsto1-e^{-u}$ is $1$-Lipschitz, so with
$g(u)=1-e^{-E_1(u)}$,
\[
\Bigl|\sum_{m\ge1}\Bigl(1-e^{-\sum_{j>m}f(j)}\Bigr)-\sum_{m\ge1}g\bigl(\tfrac mt\bigr)\Bigr|\le\sum_{m\ge1}f(m)=-\log(1-e^{-1/t})=O(\log t).
\]
The term $m=0$ is at most $1$. Finally $g$ is decreasing, integrable ($g\le1$, and $g(u)\le E_1(u)\le e^{-u}$ for $u\ge1$), so
$\frac1t\sum_{m\ge1}g(\frac mt)\to\int_0^\infty g$, a Riemann sum (Tome~4, Chapter~4). So the right-hand side is
$t\int_0^\infty g+o(t)$.

(d) If $\mathbb E[L_n]=\lambda n+o(n)$, then $\sum_n(1-x)x^n\mathbb E[L_n]=\lambda\frac x{1-x}+o\bigl(\frac1{1-x}\bigr)$ (an
Abelian argument: split the $o(n)$ at a large $n_0$), and $\frac1{1-x}\sim t$. Comparing with (c), $\lambda=\int_0^\infty g$.
Integrating by parts, with $g'(u)=-\frac{e^{-u}}ue^{-E_1(u)}$ and $ug(u)\to0$ at $0$ and $\infty$:
\[
\int_0^\infty g(u)\,du=-\int_0^\infty u\,g'(u)\,du=\int_0^\infty e^{-u-E_1(u)}\,du=0.6243299885\ldots
\]
\tools Tome~2, Chapter~4 (Cauchy's formula); Tome~3, Chapters~2 and~9 (Poisson variables); Tome~4, Chapters~4 and~5 (Riemann sums,
the exponential integral).
\leads Golomb found the constant numerically in 1964, and Shepp and Lloyd (1966) proved that $\frac{\mathbb E[L_n]}n$ converges,
with this very Poissonisation. The same $\lambda$ is the limit of the mean of $\frac{\log P^+(n)}{\log n}$, $P^+(n)$ the largest prime
factor of $n$: Dickman had found it in 1930. Figure~\ref{fig:longest} shows the convergence.
\end{solution}

\begin{solution}{pb:conjugate}
\idea Two permutations are conjugate when they have the same cycle type, so $P_n$ is a sum of squared class probabilities, and
the class probabilities factorise over cycle lengths. For the lower bound, the types with one huge cycle of length $n-k$ and a
small rest behave like a copy of the problem of size $k$, at the price $\frac1{(n-k)^2}$.

\writeup (a) The class of type $\lambda$ (with $m_j$ cycles of length $j$) has probability $\frac1{z_\lambda}$, $z_\lambda=\prod_jj^{m_j}m_j!$
(Cauchy's formula). Conjugate means same type (Tome~2, Chapter~4), so $P_n=\sum_{\lambda\vdash n}z_\lambda^{-2}$, and
\[
\sum_{n\ge0}P_nx^n=\prod_{j\ge1}\sum_{m\ge0}\frac{x^{jm}}{j^{2m}m!^2},
\]
since choosing a type is choosing $m_j$ for each $j$. All terms are positive, so at $x=1$ both sides are equal in $[0,\infty]$. The
product converges: the $j$-th factor is $1+\frac1{j^2}+O(\frac1{j^4})$. Numerically it is $4.2634\ldots$ (each factor is
$I_0(\frac2j)$, a Bessel function).

(b) Let $k<\frac n2$ and $\mu\vdash k$. The type $\lambda=\mu\cup\{n-k\}$ has a unique cycle of length $n-k$, so
$z_\lambda=(n-k)z_\mu$, and different pairs $(k,\mu)$ give different types. Keeping only these types,
\[
P_n\ge\sum_{k<n/2}\frac1{(n-k)^2}\sum_{\mu\vdash k}z_\mu^{-2}\ge\frac1{n^2}\sum_{k<n/2}P_k,
\]
and the last sum tends to $\sum_kP_k$.
\tools Tome~2, Chapter~4 (conjugacy classes of $\mathfrak S_n$); Tome~4, Chapter~3 (infinite products).
\leads In fact $P_n\sim\frac{4.2634\ldots}{n^2}$: Flajolet, Fusy, Gourdon, Panario and Pouyanne (2006) proved it with a ``hybrid'' of
Darboux's method and singularity analysis, needed because the generating function has the whole unit circle as natural boundary.
The lower bound is the main term: two random permutations are conjugate mostly when both have one giant cycle of the same length.
\end{solution}

\begin{solution}{pb:commutators}
\idea $N$ is a class function, so it is determined by its inner products with the irreducible characters. Summing $\chi$ over all
commutators $[x,y]$, the sum over $y$ of $\rho(y)\rho(x)^{-1}\rho(y)^{-1}$ averages $\rho(x)^{-1}$ into a scalar, by Schur's lemma.

\writeup (a) If $[x,y]=g$ then $[hxh^{-1},hyh^{-1}]=hgh^{-1}$, so $N$ is a class function. Let $\rho$ be irreducible with
character $\chi$ of degree $d$. For any matrix $M$, $\sum_y\rho(y)M\rho(y)^{-1}$ commutes with every $\rho(z)$, so it is a scalar
(Schur's lemma, Tome~2, Chapter~11), and comparing traces it is $\frac{|G|\operatorname{tr}M}dI$. With $M=\rho(x)^{-1}$:
\[
\sum_y\chi(xyx^{-1}y^{-1})=\operatorname{tr}\Bigl(\rho(x)\sum_y\rho(y)\rho(x)^{-1}\rho(y)^{-1}\Bigr)=\frac{|G|}d\chi(x)\chi(x^{-1})=\frac{|G|}d|\chi(x)|^2.
\]
Summing over $x$, $\sum_{x,y}\chi([x,y])=\frac{|G|^2}d$, which is real; so
$\langle N,\chi\rangle=\frac1{|G|}\sum_gN(g)\overline{\chi(g)}=\frac1{|G|}\overline{\sum_{x,y}\chi([x,y])}=\frac{|G|}d$. The irreducible characters
are an orthonormal basis of class functions (Tome~2, Theorem~11.10), so $N=\sum_\chi\frac{|G|}{\chi(1)}\chi$.

(b) The table of $A_5$ on the classes $e$, $(12)(34)$, $(123)$, $(12345)$, $(13524)$ has rows $(1,1,1,1,1)$,
$(3,-1,0,\phi,1-\phi)$, $(3,-1,0,1-\phi,\phi)$, $(4,0,1,-1,-1)$, $(5,1,-1,0,0)$, with $\phi=\frac{1+\sqrt5}2$. So
\[
\frac{N(g)}{60}=1+\frac{\chi_2(g)+\chi_3(g)}3+\frac{\chi_4(g)}4+\frac{\chi_5(g)}5
\]
is $5$ at $e$, $1-\frac23+\frac15=\frac8{15}$ at $(12)(34)$, $1+\frac14-\frac15=\frac{21}{20}$ at $(123)$, and $1+\frac13-\frac14=\frac{13}{12}$ at
both $5$-cycles (since $\phi+1-\phi=1$). So $N=300,32,63,65,65$: never $0$. Check: $300+15\cdot32+20\cdot63+24\cdot65=3600=60^2$.
\tools Tome~2, Chapter~11 (Schur's lemma, orthogonality, the table of $A_5$).
\leads $N(e)=|G|\cdot k(G)$ recovers the count of commuting pairs. Ore conjectured in 1951 that every element of every finite
non-abelian simple group is a commutator; Liebeck, O'Brien, Shalev and Tiep proved it in 2010, with this formula as the main tool:
they show that $\sum_\chi\frac{\chi(g)}{\chi(1)}\ne0$ by bounding character values.
\end{solution}

\begin{solution}{pb:frobschur}
\idea $\chi(x^2)$ is the trace of $\rho(x)^2$, which is the difference of the characters of the symmetric and alternating squares.
So $\nu(\chi)$ is the difference of the dimensions of their invariants, and $V\otimes V$ has at most one invariant line. For (b),
$r$ is a class function and its inner products with characters are the $\nu$'s.

\writeup (a) If $\rho(x)$ has eigenvalues $\lambda_1,\dots,\lambda_d$, then $\operatorname{Sym}^2V$ has eigenvalues $\lambda_i\lambda_j$,
$i\le j$, and $\Lambda^2V$ has $\lambda_i\lambda_j$, $i<j$. So their characters are $\frac{\chi(x)^2\pm\chi(x^2)}2$, and
$\chi(x^2)=\chi_S(x)-\chi_\Lambda(x)$. Hence $\nu(\chi)=\langle\chi_S,1\rangle-\langle\chi_\Lambda,1\rangle=s-a$, where $s,a\ge0$ are the
dimensions of the invariant vectors. Their sum is $\langle\chi^2,1\rangle=\frac1{|G|}\sum\chi(x)^2=\langle\chi,\bar\chi\rangle$, which is
$1$ if $\bar\chi=\chi$ and $0$ otherwise ($\bar\chi$ is also an irreducible character, of the dual representation). So
$(s,a)\in\{(0,0),(1,0),(0,1)\}$ and $\nu\in\{0,1,-1\}$, with $\nu\ne0$ exactly when $\chi$ is real-valued.

(b) $r$ is a class function, and $\langle r,\chi\rangle=\frac1{|G|}\sum_gr(g)\overline{\chi(g)}=\frac1{|G|}\sum_x\overline{\chi(x^2)}=\nu(\chi)$,
real by (a). So $r=\sum\nu(\chi)\chi$.

(c) In $Q_8$, $x^2=1$ only for $x=\pm1$, so $r(1)=2$. The characters are four of degree $1$ (real, since $Q_8/\{\pm1\}\cong(\Z/2\Z)^2$, with
$\nu=1$) and one of degree $2$. So $2=4\cdot1+2\nu$, and $\nu=-1$.

(d) If $\rho$ is real, averaging any positive definite symmetric bilinear form over $G$ gives an invariant one, a non-zero vector
of $\operatorname{Sym}^2V^*\cong\operatorname{Sym}^2V$ (as $\chi$ is real, $V\cong V^*$). So $s=1$ and $\nu=1$. For $\mathfrak S_n$ all
$\nu$ are then $1$, and evaluating (b) at $e$: $\#\{\sigma:\sigma^2=e\}=r(e)=\sum_\chi\chi(1)$.
\tools Tome~2, Chapter~11 (characters, orthogonality); Tome~0, Chapter~7 (bilinear forms, eigenvalues).
\leads For $n=4$: degrees $1,1,2,3,3$, sum $10$, and $\mathfrak S_4$ has $1+6+3=10$ involutions with the identity. Since the degrees count
standard Young tableaux, (d) says that there are as many standard Young tableaux with $n$ boxes as involutions; the
Robinson--Schensted correspondence gives a bijection. The value $\nu=-1$ marks the \emph{quaternionic} representations.
\end{solution}

\begin{solution}{pb:ncycles}
\idea The class formula expresses the count as a sum over characters of $\frac{\chi(c)^3}{\chi(1)}$, and only the exterior powers
of $V$ survive, with $\chi(c)=\pm1$. What is left is a sum of reciprocals of binomial coefficients, which a Beta integral evaluates.

\writeup (a) The eigenvalues of $c$ on $\C^n$ are the $n$-th roots of unity $\zeta^j$, $0\le j<n$ (the eigenvectors are
$(\zeta^{jk})_k$), and $V$ is the sum of the eigenspaces with $j\ne0$. So $\chi_r(c)=e_r(\zeta,\dots,\zeta^{n-1})$, the elementary
symmetric function, and
\[
\sum_r\chi_r(c)t^r=\prod_{j=1}^{n-1}(1+t\zeta^j)=\frac{1-(-t)^n}{1+t}=\sum_{r=0}^{n-1}(-t)^r,
\]
using $\prod_{j=0}^{n-1}(1+t\zeta^j)=1-(-t)^n$ (substitute $X=-\frac1t$ in $\prod(X-\zeta^j)=X^n-1$). So $\chi_r(c)=(-1)^r$. The
centraliser of $c$ is the cyclic group it generates, of order $n$, and by column orthogonality (Tome~2, Chapter~11)
$\sum_\chi|\chi(c)|^2=n$. The $n$ characters $\chi_r$ already contribute $n$, so all other irreducible characters vanish at $c$.

(b) For irreducible $\rho$ of degree $d$ and a class $A$, the matrix $\sum_{a\in A}\rho(a)$ commutes with $\rho(G)$, so it is the
scalar $\frac{|A|\chi(a)}d$ (compare traces). Then
\[
\#\{(a,b)\in A\times B:ab=g\}=\frac1{|G|}\sum_\chi\chi(1)\sum_{a,b}\chi\bigl(ab\,g^{-1}\bigr),
\]
because $\frac1{|G|}\sum_\chi\chi(1)\chi(h)$ is $1$ if $h=e$ and $0$ otherwise (the regular character). And
$\sum_{a,b}\chi(abg^{-1})=\operatorname{tr}\bigl(\sum_a\rho(a)\sum_b\rho(b)\rho(g)^{-1}\bigr)=\frac{|A|\chi(a)}d\frac{|B|\chi(b)}d\chi(g^{-1})$. Since
$\chi(g^{-1})=\overline{\chi(g)}$, the formula follows.

(c) $\frac1{\binom mr}=(m+1)\int_0^1x^r(1-x)^{m-r}dx$ (Beta integral, Tome~4, Chapter~5). Summing the geometric progression,
$\sum_{r=0}^m(-x)^r(1-x)^{m-r}=(1-x)^{m+1}-(-x)^{m+1}$, since the ratio is $\frac{-x}{1-x}$ and $(1-x)-(-x)=1$. Integrating,
$\frac1{m+2}+\frac{(-1)^m}{m+2}$; multiply by $m+1$.

(d) Take $A=B$ the class of $c$, of size $(n-1)!$, and $g=c$. Only the $\chi_r$ contribute, with $\chi_r(c)^3=(-1)^r$ and
$\chi_r(1)=\binom{n-1}r$:
\[
\frac{(n-1)!^2}{n!}\sum_{r=0}^{n-1}\frac{(-1)^r}{\binom{n-1}r}=\frac{(n-1)!}n\cdot\frac n{n+1}\bigl(1+(-1)^{n-1}\bigr),
\]
which is $\frac{2(n-1)!}{n+1}$ for $n$ odd and $0$ for $n$ even.
\tools Tome~2, Chapter~11 (characters, Schur's lemma, orthogonality); Tome~4, Chapter~5 (Beta function); Tome~0, Chapter~7 (eigenvalues).
\leads For $n$ even the answer is obvious: an $n$-cycle is odd and a product of two odd permutations is even. For $n=5$ there are $8$
factorisations, for $n=7$, $180$. Stanley (1981) and Zagier (1995) studied the number of cycles of the product of two random
permutations from given classes, a question that is now part of the theory of maps on surfaces.
\end{solution}

\begin{solution}{pb:denes}
\idea Follow the cycles as the transpositions are multiplied one at a time. Starting from $n$ cycles, each step changes the number
of cycles by $\pm1$, so reaching one cycle in $n-1$ steps requires a merge every time, and merges are exactly the edges that join two
components of the graph built so far.

\writeup \emph{The rule.} Let $\sigma$ be a permutation and $a\ne b$. If $a$ and $b$ are in different cycles
$(a\,x_1\,\dots\,x_r)$ and $(b\,y_1\,\dots\,y_s)$, then $\sigma(a\,b)$ has the cycle $(a\,y_1\,\dots\,y_s\,b\,x_1\,\dots\,x_r)$ instead (check:
$a\mapsto b\mapsto y_1$, \dots, $b\mapsto a\mapsto x_1$): a merge. If they are in the same cycle, the same computation read backwards is a
split. So the number of cycles changes by $\pm1$.

(a) Let $\pi_i=\tau_1\cdots\tau_i$, $\pi_0=e$, with $n$ cycles. $\pi_{n-1}$ has one cycle if and only if every step merges. Let $T_i$ be
the graph with edges $\{a_1,b_1\},\dots,\{a_i,b_i\}$. If every step up to $i$ merges, then by induction the cycles of $\pi_i$ have the
same supports as the components of $T_i$: a merge joins the cycles of $a_{i+1}$ and $b_{i+1}$, and the new edge joins their
components. So step $i+1$ merges if and only if $a_{i+1}$ and $b_{i+1}$ are in different components of $T_i$, that is, if the new
edge creates no cycle. All $n-1$ steps merge if and only if $T_{n-1}$ is acyclic with $n-1$ edges, that is, a tree.

(b) By (a), the pairs (tree, ordering of its edges) are in bijection with the factorisations of $n$-cycles into $n-1$
transpositions: there are $n^{n-2}(n-1)!$ of them, by Cayley's formula. Conjugating by a permutation $\pi$ maps factorisations of
$c$ bijectively onto those of $\pi c\pi^{-1}$, and all $(n-1)!$ $n$-cycles are conjugate. So each has $\frac{n^{n-2}(n-1)!}{(n-1)!}=n^{n-2}$
factorisations.
\tools Tome~2, Chapter~4 (cycles, conjugacy); Tome~3, Chapter~12 (Cayley's formula).
\leads Dénes (1959). Counting factorisations of permutations into transpositions gives the \emph{Hurwitz numbers}, which count
branched covers of the sphere; Goulden and Jackson found a formula for every cycle type in 1997. Figure~\ref{fig:tree} shows one
factorisation.
\end{solution}

\section*{Map of the chapter}

\begin{center}
\small
\adjustbox{max width=\linewidth}{\begin{tabular}{lll}
Problem & Topic & Tomes used\\\hline
\ref{pb:dice} & dice up to rotation & 2\\
\ref{pb:patterns} & colourings up to colours (de Bruijn) & 2\\
\ref{pb:necklacering} & necklaces, cyclotomic identity & 1, 4\\
\ref{pb:unlabelled} & almost all graphs are asymmetric & 2, 4\\
\ref{pb:commuting} & commuting tuples of permutations & 2, 4, 1\\
\ref{pb:landau} & groups with few classes & 2, 1\\
\ref{pb:rank} & fixed points and rank & 2, 3\\
\ref{pb:derangements} & Jordan, Cameron--Cohen & 2, 3\\
\ref{pb:glfix} & fixed vectors of a random matrix & 0, 2, 3\\
\ref{pb:ksubsets} & characters on $k$-subsets & 2, 0\\
\ref{pb:samecycle} & points in the same cycle & 3, 0\\
\ref{pb:foata} & cycles and records & 2, 3\\
\ref{pb:roots} & solutions of $\sigma^k=e$ & 4, 1, 2\\
\ref{pb:oddcycles} & odd and even cycles & 4, 3\\
\ref{pb:squares} & square permutations & 2, 4\\
\ref{pb:poissoncycles} & Poisson cycle counts & 3, 5\\
\ref{pb:golomb} & longest cycle, Golomb--Dickman & 2, 3, 4\\
\ref{pb:conjugate} & two conjugate permutations & 2, 4\\
\ref{pb:commutators} & Frobenius's commutator formula & 2\\
\ref{pb:frobschur} & Frobenius--Schur, square roots & 2, 0\\
\ref{pb:ncycles} & products of two $n$-cycles & 2, 4, 0\\
\ref{pb:denes} & trees and transpositions & 2\\
\end{tabular}}
\end{center}

\section*{Sources and further reading}

\begin{itemize}
\item H.~Weyl, \emph{Symmetry}, Princeton University Press, 1952, p.~5 (the epigraph).
\item A.~Cayley, On the analytical forms called trees, with application to the theory of chemical combinations,
\emph{Report of the British Association for the Advancement of Science} (1875), 257--305; N.~L.~Biggs, E.~K.~Lloyd and
R.~J.~Wilson, \emph{Graph Theory 1736--1936}, Oxford University Press, 1976, Chapter~4 (Cayley's counts and their history).
\item G.~Pólya, Kombinatorische Anzahlbestimmungen für Gruppen, Graphen und chemische Verbindungen, \emph{Acta Mathematica} 68 (1937),
145--254; English translation in G.~Pólya and R.~C.~Read, \emph{Combinatorial Enumeration of Groups, Graphs, and Chemical
Compounds}, Springer, 1987.
\item P.~M.~Neumann, A lemma that is not Burnside's, \emph{The Mathematical Scientist} 4 (1979), 133--141.
\item N.~G.~de Bruijn, Generalization of Pólya's fundamental theorem in enumerative combinatorial analysis, \emph{Indagationes
Mathematicae} 21 (1959), 59--69.
\item N.~Metropolis and G.-C.~Rota, Witt vectors and the algebra of necklaces, \emph{Advances in Mathematics} 50 (1983), 95--125.
\item F.~Harary and E.~M.~Palmer, \emph{Graphical Enumeration}, Academic Press, 1973 (unlabelled graphs, Oberschelp's expansion).
\item J.~Bryan and J.~Fulman, Orbifold Euler characteristics and the number of commuting $m$-tuples in the symmetric groups,
\emph{Annals of Combinatorics} 2 (1998), 1--6.
\item P.~J.~Cameron and A.~M.~Cohen, On the number of fixed point free elements in a permutation group, \emph{Discrete
Mathematics} 106/107 (1992), 135--138; J.-P.~Serre, On a theorem of Jordan, \emph{Bulletin of the American Mathematical Society}
40 (2003), 429--440.
\item J.~Fulman, Random matrix theory over finite fields, \emph{Bulletin of the American Mathematical Society} 39 (2002), 51--85.
\item D.~Livingstone and A.~Wagner, Transitivity of finite permutation groups on unordered sets, \emph{Mathematische Zeitschrift}
90 (1965), 393--403.
\item R.~Arratia, A.~D.~Barbour and S.~Tavaré, \emph{Logarithmic Combinatorial Structures: a Probabilistic Approach}, EMS, 2003.
\item L.~A.~Shepp and S.~P.~Lloyd, Ordered cycle lengths in a random permutation, \emph{Transactions of the American Mathematical
Society} 121 (1966), 340--357.
\item J.~Blum, Enumeration of the square permutations in $S_n$, \emph{Journal of Combinatorial Theory A} 17 (1974), 156--161.
\item P.~Flajolet, É.~Fusy, X.~Gourdon, D.~Panario and N.~Pouyanne, A hybrid of Darboux's method and singularity analysis in
combinatorial asymptotics, \emph{Electronic Journal of Combinatorics} 13 (2006), R103.
\item M.~W.~Liebeck, E.~A.~O'Brien, A.~Shalev and P.~H.~Tiep, The Ore conjecture, \emph{Journal of the European Mathematical
Society} 12 (2010), 939--1008.
\item W.~Fulton and J.~Harris, \emph{Representation Theory: A First Course}, Springer, 1991 (exterior powers of the standard
representation, Frobenius--Schur indicators).
\item R.~P.~Stanley, Factorization of permutations into $n$-cycles, \emph{Discrete Mathematics} 37 (1981), 255--262.
\item J.~Dénes, The representation of a permutation as the product of a minimal number of transpositions, and its connection
with the theory of graphs, \emph{Publications of the Mathematical Institute of the Hungarian Academy of Sciences} 4 (1959), 63--71.
\end{itemize}
